\documentclass[11pt]{article}
\usepackage[a4paper,margin=0.65in]{geometry}
\usepackage[T1]{fontenc}
\usepackage{amsmath,amssymb,amsfonts}
\usepackage{graphicx}
\usepackage{pdfpages}
\usepackage[english,shorthands=off]{babel}
\usepackage{fancyhdr}
\usepackage{array}
\usepackage{rotating}
\usepackage[bookmarks=false,colorlinks=false]{hyperref}
\usepackage[protrusion=true,expansion=false]{microtype}
\emergencystretch=3em
\tolerance=600
\providecommand{\modd}{\mathrm{mod}}
\providecommand{\Disc}{\mathrm{Disc}}
\providecommand{\canont}{}
\providecommand{\asterism}{*}
\providecommand{\Hessian}{H}
\providecommand{\tome}{}
\providecommand{\page}{p.}
\pagestyle{fancy}\fancyhf{}\fancyhead[L]{Pages 138--162}\fancyhead[R]{Cayley --- Collected Papers, Vol.\ IX}\renewcommand{\headrulewidth}{0.4pt}
\setlength{\headheight}{15pt}

\begin{document}

\begin{center}
{\Large\bfseries A Memoir on the Transformation of Elliptic Functions}\\[4pt]
{\small (Continuation of paper 578 --- pp.\ 118--142 of Vol.\ IX)}
\end{center}

\section*{[p.~118] Tables for the cases $n=7$ and $n=11$ \emph{(continued)}}

\textit{[Table: the second equation-system for the septic case $n=7$. The
rows are labelled by powers $\Omega^k$ ($k=0,\ldots,8$) and the columns by
$k^9, k^8, k^7, k^6, k^5, k^4, k$, $1$. The non-zero entries are arranged so
that the equation reads, on summation by rows,]}
\[
\begin{aligned}
\Omega^0 &: \ +1, \\
\Omega^1 &: \ -64\,k^8 + 56\,k^6, \\
\Omega^2 &: \ -112\,k^7 + 140\,k^5, \\
\Omega^3 &: \ -112\,k^6 + 56\,k^4, \\
\Omega^4 &: \ +70\,k^5, \\
\Omega^5 &: \ +56\,k^4 - 112\,k^2, \\
\Omega^6 &: \ +140\,k^3 - 112\,k, \\
\Omega^7 &: \ +56\,k^2 - 64, \\
\Omega^8 &: \ +1,
\end{aligned}
\]
the column-sums of which (read in the order $k^9,k^8,k^7,k^6,k^5,k^4,k,1$)
are
\[
-64,\ -112,\ 0,\ +352,\ 0,\ -112,\ -64,\ =71\cdot 11;
\]
that is, $\Omega=1$, we have
\[
-16(k-1)^2(4k^2+3k+1)(k^2+3k+4)=0.
\]

\vspace{1ex}
\noindent\textit{[Table: the equation-system for $n=11$, endecadic
transformation. Rows are labelled $\Omega^0,\ldots,\Omega^{12}$ and columns
by $k^{10}, k^9, k^8, k^7, k^6, k^5, k^4, k^3, k^2, k, 1$. The non-zero
entries, row by row, are:]}
\[
\begin{aligned}
\Omega^0 &:\ +1, \\
\Omega^1 &:\ -1024\,k^{10} + 1408\,k^9 - 396\,k^8 - 12, \\
\Omega^2 &:\ -5632\,k^9 + 4400\,k^8 + 1298\,k^7 + 66, \\
\Omega^3 &:\ -16192\,k^8 + 16368\,k^7 - 396\,k^5 - 220, \\
\Omega^4 &:\ -18656\,k^7 + 19151\,k^6 + 495, \\
\Omega^5 &:\ -16016\,k^6 - 1144\,k^5 + 16368\,k^4 - 792, \\
\Omega^6 &:\ +4400\,k^5 - 7876\,k^4 + 4400\,k^3 + 924, \\
\Omega^7 &:\ +16368\,k^4 - 1144\,k^3 - 16016\,k^2 - 792, \\
\Omega^8 &:\ +19151\,k^3 - 18656\,k^2 + 495, \\
\Omega^9 &:\ -396\,k^4 + 16368\,k^2 - 16192\,k - 220, \\
\Omega^{10} &:\ +1298\,k^3 + 4400\,k - 5632 + 66, \\
\Omega^{11} &:\ -396\,k^2 + 1408 - 1024 - 12, \\
\Omega^{12} &:\ +1,
\end{aligned}
\]
the column-sums of which (read across the columns
$k^{10},k^9,\ldots,k,1$) are
\[
\begin{aligned}
&-1024,\ -32208,\ -18656,\ -1936,\ -7876,\ -1936,\ -18656,\ -32208,\ -1024,\\
&+1408,\ +8800,\ +32736,\ +40900,\ +32736,\ +8800,\ +1408,\\
&-5632,\ -30800,\ -9856,\ +30800,\ +33024,\ +30800,\ -9856,\ -30800,\ -5632,\\
&-1024 + 94624.
\end{aligned}
\]

\section*{[p.~119] Equation-systems for the cases $n=3,5,7,9,11$. Art.\ Nos.\ 8 to 10.}

\textbf{8.} $n=3$, cubic transformation. \quad $k=u^4$, \
$\Omega=\dfrac{u^9}{v}$ (here and in the other cases).
\[
P=\alpha,\ \ Q=\beta. \quad \text{The condition here is}
\]
\[
k^2\alpha^2\alpha^2 + (2\alpha\beta+\beta^2)^2 = \Omega\bigl[(\alpha^2+2\alpha\beta)^2 + \beta^2 \alpha^2\bigr],
\]
and the system of equations thus is
\[
k\alpha^2 = \Omega\beta^2,\qquad
2\alpha\beta + \beta^2 = \Omega k\,(\alpha^2 + 2\alpha\beta),
\]
and similarly in the other cases; for these it will be enough to write down
the equation-systems.

\medskip
$n=5$, quintic transformation.
\[
P=\alpha+\gamma x^2,\ \ Q=\beta.
\]
\[
k^2\alpha^2 = \Omega\gamma^2,
\]
\[
2\alpha\gamma + 2\alpha\beta + \beta^2 = \Omega(2\beta\gamma + 2\beta\gamma + \beta^2),
\]
\[
\gamma^2 + 2\beta\gamma = \Omega k^2(\alpha^2 + 2\alpha\beta).
\]

\medskip
$n=7$, septic transformation.
\[
P=\alpha+\gamma x^2,\ \ Q=\beta+\delta x^2.
\]
\[
k^2\alpha^2 = \Omega\delta^2,
\]
\[
k\,(2\alpha\gamma + 2\alpha\beta + \beta^2) = \Omega(\gamma^2 + 2\gamma\delta + 2\beta\delta),
\]
\[
\gamma^2 + 2\beta\gamma + 2\alpha\delta + 2\beta\delta
= \Omega k\,(2\alpha\gamma + 2\beta\gamma + 2\alpha\delta + \beta^2),
\]
\[
\delta^2 + 2\gamma\delta = \Omega k^2(\alpha^2 + 2\alpha\beta).
\]

\medskip
$n=9$, enneadic transformation.
\[
P=\alpha+\gamma x^2 + \epsilon x^4,\ \ Q=\beta+\delta x^2.
\]
\[
k^2\alpha^2 = \Omega\epsilon^2,
\]
\[
k\,(2\alpha\gamma + 2\alpha\beta + \beta^2)
= \Omega(2\gamma\epsilon + 2\alpha\delta + \delta^2),
\]
\[
2\alpha\epsilon + \gamma^2 + 2\alpha\delta + 2\gamma\beta + 2\beta\delta
= \Omega(2\alpha\epsilon + \gamma^2 + 2\gamma\delta + 2\delta\beta + 2\beta\delta),
\]
\[
2\gamma\epsilon + 2\gamma\delta + 2\epsilon\beta + \delta^2
= \Omega k\,(2\beta\gamma + 2\alpha\delta + 2\gamma\beta + \beta^2),
\]
\[
\epsilon^2 + 2\delta\epsilon = \Omega k^2(\alpha^2 + 2\alpha\beta).
\]

\medskip
$n=11$, endecadic transformation.
\[
P=\alpha+\gamma x^2 + \epsilon x^4,\ \ Q=\beta+\delta x^2 + \zeta x^4.
\]
\[
k^2\alpha^2 = \Omega\zeta^2,
\]
\[
k\,(2\alpha\gamma + 2\alpha\beta + \beta^2)
= \Omega(\epsilon^2 + 2\epsilon\zeta + 2\delta\zeta),
\]
\[
k\,(2\alpha\epsilon + \gamma^2 + 2\alpha\delta + 2\gamma\beta + 2\beta\delta)
= \Omega(2\gamma\epsilon + 2\gamma\zeta + 2\epsilon\delta + 2\beta\zeta + \delta^2),
\]
\[
2\gamma\epsilon + 2\alpha\zeta + 2\gamma\delta + 2\epsilon\beta + 2\beta\zeta + \delta^2
= \Omega k\,(2\alpha\epsilon + \gamma^2 + 2\alpha\zeta + 2\gamma\delta + 2\epsilon\beta + 2\beta\delta),
\]
\[
\epsilon^2 + 2\gamma\zeta + 2\epsilon\delta + 2\delta\zeta
= \Omega k^2(2\alpha\gamma + 2\alpha\delta + 2\gamma\beta + \beta^2),
\]
\[
2\epsilon\zeta + \zeta^2 = \Omega k^2(\alpha^2 + 2\alpha\beta).
\]

And so on.

\section*{[p.~120]}

\textbf{9.} It will be noticed that if the coefficients of $P+Qx$ taken in
order are
\[
\alpha,\ \beta,\ \gamma,\ \ldots,\ \sigma,\ \tau,
\]
then in every case the first and last equations are
\[
k\alpha^2 = \Omega\tau^2,\qquad
\tau^2 + 2\rho\sigma\tau = \Omega^{\frac{1}{2}(n-1)} k^{\frac{1}{2}(n-1)}\!\bigl(\alpha^2 + 2\alpha\beta\bigr).
\]
Putting in the first of these $k=u^4$, $\Omega=\dfrac{u^n}{v}$, the
equation becomes
\[
\frac{u^4\alpha^2}{v} \;=\; \frac{u^n}{v}\,\tau^2,
\]
where each side is a perfect square; and in extracting the square root we
may without loss of generality take the roots positive, and write
$\tau = u^{\frac{1}{2}(n-4)}\alpha = v\alpha$.

This speciality, although it renders it proper to employ ultimately
$u, v$ in place of $k, \Omega$, produces really no depression of order
(viz.\ the $\Omega k$-form of the modular equation is found to be of the
same order in $\Omega$ that the standard or $uv$-form is in $v$), and is
in another point of view a disadvantage, as destroying the uniformity of
the several equations: in the discussion of order I consequently retain
$\Omega, k$. Ultimately these are to be replaced by $u, v$; the change in
the equation-systems is so easily made that it is not necessary here to
write them down in the new form in $u, v$.

\medskip
\textbf{10.} The case $\alpha=0$ has to be considered in the discussion of
order, but we have thus only solutions which are to be rejected; in the
proper solutions $\alpha$ is not $=0$, and it may therefore for convenience
be taken to be $=1$. We have then $\sigma = u^{n-4} \div v$.

The last equation becomes therefore
\[
\tau^2 + 2\rho\sigma\tau \;=\; \Omega^{\frac{1}{2}(n-1)}k^{\frac{1}{2}(n-1)}\,(1+2\beta),
\]
or recollecting that $\beta$ is connected with the multiplier $M$ by the
relation $\dfrac{1}{M} = 1 + 2\beta$, that is,
\[
2\beta = \frac{1}{M} - 1,
\]
and substituting for $1+2\beta$ its value, the equation becomes
\[
2\rho = u^{4(n-1)}\!\left(\frac{1}{M} - \frac{u^n}{v^n}\right);
\]
that is, the first and the last coefficients are $1,\ \dfrac{u^{n-4}}{v}$,
and the second and the penultimate coefficients are each expressed in
terms of $v, M$. The cases $n=3,\ n=5$ are so far peculiar, that the only
coefficients are $\alpha,\beta$, or $\alpha,\beta,\gamma$; in the next
case $n=7$, the only coefficients are $\alpha,\beta,\gamma,\delta$, and we
have in this case all the coefficients expressed as above.

\section*{[p.~121] The $\Omega k$-form --- Order of the Systems. Art.\ Nos.\ 11 to 22.}

\textbf{11.} In the general case, $n$ an odd number, we have $\Omega$, and
$\tfrac{1}{2}(n+1)$ coefficients connected by a system of $\tfrac{1}{2}(n+1)$
equations of the form
\[
\Omega = \frac{U}{U'} = \frac{V}{V'} = \cdots,
\]
where $U, V,\ldots, U', V',\ldots$ are given quadric functions of the
coefficients. Omitting the $\Omega$ ($=$), there remains a system of
$\tfrac{1}{2}(n-1)$ equations of the form $\dfrac{U}{U'} = \dfrac{V}{V'} = \cdots$,
or say
\[
\left|\begin{array}{ccc} U, & V, & W,\ldots \\ U', & V', & W',\ldots \end{array}\right| = 0,
\]
which determine the ratios $\alpha:\beta:\gamma:\ldots$ of the coefficients;
and to each set of ratios there corresponds a single value of $\Omega$. The
order of the system, or number of sets of ratios, is
$=\tfrac{1}{2}(n+1)\cdot 2^{\frac{1}{2}(n-1)}$,
$\;=\tfrac{1}{2}(n+1)\cdot 2^{\frac{1}{2}(n-1)}$;
and this is consequently the number of values of $\Omega$, or the order
of the equation for the determination of $\Omega$; viz.\ but for reduction,
the order in $\Omega$ of the $\Omega k$-modular equation would be
$= (n+1)\cdot 2^{\frac{1}{2}(n-1)}$. In the case $n=3$, this is $4$, which
is right, but for any larger value of $n$ the order is far too high; in
fact, assuming (as the case is) that the order is equal to the order in
$v$ of the $uv$-form, the order should for a prime value of $n$ be
$=n+1$, and for a composite value not containing any square factor be
$=$ the sum of the divisors of $n$. I do not attempt a general
investigation, but confine myself to showing in what manner the
reductions arise.

\medskip
\textbf{12.} I will first consider the cubic transformation; here, writing
for convenience $\dfrac{\alpha}{\beta} = \theta$, the equations give
\[
\frac{k\theta^2}{2\theta+1} \;=\; \frac{1}{k(\theta^2+2\theta)},
\quad\text{that is,}\quad k^2\theta^2(\theta+2) - (2\theta+1) = 0,
\]
and
\[
k\theta^2 = \Omega.
\]
The equation in $\theta$ gives $(k^2\theta^2 - 1)^2 - 4\theta\,(k^2\theta^2 - 1)^2 = 0$,
and we have thence
\[
k(\Omega^2 - 1)^2 \;=\; 4\Omega\,(k\Omega - 1)^2,
\]
that is, $k(\Omega^2-1)^2 - 4\Omega(k\Omega-1)^2 = 0$,
\[
k\Omega^4 - 4k\Omega^3 + 6k\Omega^2 - 4\Omega + k = 0,
\]
the modular equation; and then $-1 + k^2\theta^2 \mid 2\theta\,(k^2\theta^2 - 1) = 0$,
that is, $\beta = 1 + 2\theta(k\Omega - 1) = 0$, or $2\theta = -\dfrac{1}{k\Omega - 1}$;
which is to say we have $\alpha=\beta=-1$, $\beta = 2(1-k\Omega)$; consequently
\[
\lambda = \Omega k,\quad\text{and}\quad
\frac{1-y}{1+y} \;=\; \frac{1-x}{1+x}\!\left[\frac{P - Q x}{P + Q x}\right]
\;=\; \frac{1-x}{1+x}\!\left[\frac{-1 + 2(k\Omega - 1)x}{-1 - 2(k\Omega - 1)x}\right],
\]
which completes the theory.

\section*{[p.~122]}

\textbf{13.} Reproducing for this case the general theory, it appears \emph{\`a
priori} that $\Omega$ is determined by a quartic equation; in fact, from
the original equations eliminating $\Omega$, we have an equation
\[
\left|\begin{array}{cc} U, & V \\ U', & V' \end{array}\right| = 0,
\]
where $U, U', V, V'$ are quartic functions of $\alpha, \beta$; that is, the
ratio $\alpha:\beta$ has four values, and to each of these there corresponds
a single value of $\Omega$; viz.\ $\Omega$ is determined by a quartic equation.

\medskip
\textbf{14.} Considering next the case $n=5$, the quintic transformation;
the elimination of $\Omega$ gives the equations
\[
\frac{U}{U'} = \frac{V}{V'} = \frac{W}{W'},
\]
where $U, U'$, \&c.\ are all quadric functions of $\alpha,\beta,\gamma$. We
have thence $4\cdot 4 - 2\cdot 2 = 12$ sets of values of $\alpha:\beta:\gamma$;
viz.\ considering $\alpha,\beta,\gamma$ as coordinates \emph{in piano}, the
curves $UV' - U'V = 0$, $UW' - U'W = 0$ are quartic curves intersecting in
$16$ points; but among these are included the four points $\alpha=0,\ U=0$
(in fact, the point $\alpha=0,\ U'=0$ four times), which are not points of
the curve $VW' - V'W = 0$; there remain therefore $16-4,\ =12$
intersections, agreeing with the general value $(n+1)\cdot 2^{\frac{1}{2}(n-1)}$.
Hence $\Omega$ is in the first instance determined by an equation of the
order $12$; but the proper order being $=6$, there must be a factor of the
order $6$ to be rejected. To explain this and to determine the factor,
observe that the equations in question are
\[
k^2\alpha^2(2\beta\gamma + 2\beta\gamma + \beta^2) - \gamma^2(2\alpha\gamma + 2\alpha\beta + \beta^2) = 0,
\]
\[
k^2\alpha^2(\alpha+2\beta) - \gamma^2(\gamma + 2\beta) = 0;
\]
at the point $\alpha=0,\ \gamma=0$, the first of these has a double point,
the second a triple point; or there are at the point in question $6$
intersections; but $4$ of these are the points which give the foregoing
reduction $16-4=12$; we have thus the point $\alpha=0,\ \gamma=0$, counting
twice among the twelve points. Writing in the two equations $\beta=0$, the
equations become $k^2\alpha^2\gamma^2 - \alpha^2\gamma^2 = 0$,
$k^2\alpha^4 - \gamma^4 = 0$, viz.\ these will be satisfied if
$k^2\alpha^2 - \gamma^2 = 0$, that is, the curves pass through each of the
two points $(\beta=0,\ \gamma=\pm k\alpha)$, and these values satisfy (as
in fact they should) the third equation
\[
k^2(2\alpha\gamma + 2\alpha\beta + \beta^2)\,\alpha\,(\alpha+2\beta) - \gamma(\gamma+2\beta)\,(2\alpha\gamma + 2\beta\gamma + \beta^2) = 0.
\]
It is moreover easily shown that the three curves have at each of the
points in question a common tangent; viz.\ taking $A,B,C$ as current
coordinates, the tangent at the point $(\alpha,\beta,\gamma)$ of the second
curve has for its equation
\[
A\,(2\alpha^3 + 3\alpha^2\beta)\,k^2 + B\,(k^2\alpha^3 - \gamma^3) + C\,(-2\gamma^3 - 3\gamma^2\beta) = 0;
\]
and for $\beta=0,\ \gamma=\pm k\alpha$, this becomes
$2k\,A + B(k+1) + 2C = 0$, viz.\ this is the line from the point
$(\beta=0,\ \gamma=\pm k\alpha)$ to the point $(1,-2,1)$. And similarly for
the other two curves we find the same equation for the tangent.

\section*{[p.~123]}

Hence among the $12$ points are included the point $(\gamma=0,\ \alpha=0)$
twice, and the points $(\beta=0,\ \gamma=\pm k\alpha)$ each twice: we have
thus a reduction $=6$.

\medskip
\textbf{15.} Writing in the equations $\gamma=0,\ \alpha=0$, the first and
third are satisfied identically, and the second becomes
$\beta^2 = \Omega\beta^2$, that is, the equations give $\Omega=1$; writing
$\beta=0$, they become
\[
k^2\alpha^2 = \Omega\gamma^2,\quad
\alpha\gamma = \Omega\alpha\gamma,\quad
\gamma^2 = \Omega k^2 \alpha^2,
\]
viz.\ putting herein $\gamma^2 = k^2\alpha^2$, the equations again give
$\Omega=1$; hence the factor of the order $6$ is $(\Omega - 1)^6$, and the
equation of the twelfth order for the determination of $\Omega$ is
\[
(\Omega - 1)^6 \bigl\{(\Omega,1)^6\bigr\} = 0,
\]
where $(\Omega, 1)^6 = 0$ is the $\Omega k$-modular equation above written down.

\medskip
\textbf{16.} Reverting to the equation
\[
\frac{1-y}{1+y} \;=\; \frac{1-x}{1+x}\!\left(\frac{P - Q x}{P + Q x}\right)^{\!2},
\]
it is to be observed that for $\alpha=0,\ \gamma=0$, that is, $P=0$, this
becomes simply $y=x$, which is the transformation of the order $1$; the
corresponding value of the modulus $\lambda$ is $\lambda=k$, and the
equation $\lambda = \Omega^2 k$ then gives $\Omega^2 = 1$, which is
replaced by $\Omega - 1 = 0$.

If in the same equation we write $\beta=0$, that is, $Q=0$, then (without
any use of the equation $\gamma^2 = k^2\alpha^2$) we have $y=x$, the
transformation of the order $1$; but although this is so, the fundamental
equation
\[
(P^2 + 2PQ x^2 + Q^2 x^4)^2 \;=\; \Omega^2 k\,(P^4 + 2PQ + Q^2 x^4),
\]
which, putting therein $Q=0$, becomes $(P^2)^2 = \Omega^2 k^{-1} P^4$, that
is, $(k^2\alpha + \gamma)^2 = \Omega^2 k^{-1}(\alpha + \gamma x)^2$, is not
satisfied by the single relation $\Omega - 1 = 0$, but necessitates the
further relation
\[
\gamma^2 = k^2\alpha^2.
\]

The thing to be observed is that the extraneous factor $(\Omega - 1)^2$,
equated to zero, gives for $\Omega$ the value $\Omega = 1$ corresponding to
the transformation $y = x$ of the order $1$.

\medskip
\textbf{17.} Considering next $n=7$, the septic transformation; we have
here between $\alpha,\beta,\gamma,\delta$ a fourfold relation of the form
\[
\left|\begin{array}{cccc} U, & V, & W, & Z \\ U', & V', & W', & Z' \end{array}\right| = 0,
\]
where, as before, $U, V$, \&c.\ are quadric functions, and the number of
solutions is here $= 8\cdot 2^3,\ = 32$; to each of these corresponds a
single value of $\Omega$, or $\Omega$ is in the first instance determined
by an equation of the order $32$. But the order of the modular equation is
$=8$; or representing this by $\{(\Omega, 1)^8\} = 0$, the equation must be
\[
(\Omega - 1)^{24}\,\bigl\{(\Omega, 1)^8\bigr\} = 0,
\]
viz.\ there must be a special factor of the order $24$.

\section*{[p.~124]}

\textbf{18.} One way of satisfying the equations is to write therein
$\alpha=0,\ \delta=0$; the equations thus become
\[
k\beta^2 = \Omega\gamma^2,\qquad
\gamma^2 + 2\beta\gamma = \Omega k\,(2\beta\gamma + \beta^2);
\]
or putting $\beta=\alpha'$, $\gamma=\beta'$,
\[
k\alpha'^2 = \Omega\beta'^2,\qquad
\beta'^2 + 2\alpha'\beta' = \Omega k\,(\alpha'^2 + \alpha'\beta'),
\]
which (with $\alpha',\beta'$ instead of $\alpha,\beta$) are the very
equations which belong to the cubic transformation; hence a factor is
$\{(\Omega, 1)^4\}$.

Observe that for the values in question $\alpha=0,\ \delta=0$,
$P = \gamma x^2,\ Q = \beta,\ =\alpha',$
and therefore $P,Q = (\beta - Q x)^2 = \alpha'^2(\beta' - \beta' x)^2,
=\alpha'^2(P' - Q' x)^2$, if $P' = \alpha',\ Q' = \beta'$,
\[
\frac{1-y}{1+y} \;=\; \frac{1-x}{1+x}\!\left(\frac{P' - Q' x}{P' + Q' x}\right)^{\!2},
\]
which is the formula for a cubic transformation.

\medskip
\textbf{19.} The equations may also be satisfied by writing therein
$\gamma=k\alpha,\ \delta=k\beta$; in fact, substituting these values, they
become
\[
k^2\alpha^2 = \Omega k^2\beta^2,
\]
\[
2k^2\alpha^2 + k(2\alpha\beta+\beta^2) = \Omega k^2(\alpha^2+2\alpha\beta) + 2\Omega k\beta^2,
\]
\[
k^2\alpha^2 + 2k(\beta^2 + 2\alpha\beta) = 2\Omega k^2(\beta^2+2\alpha\beta) + \Omega k\beta^2,
\]
\[
k\,(\beta^2+2\alpha\beta) = \Omega k\,(\alpha^2+2\alpha\beta);
\]
the first and last of these are
\[
k\alpha^2 = \Omega\beta^2,\qquad
\beta^2 + 2\alpha\beta = \Omega k\,(\alpha^2 + 2\alpha\beta),
\]
which being satisfied the second and third equations are satisfied
identically; and these are the formulae for a cubic transformation; that
is, we again have the factor $\{(\Omega, 1)^4\}$.

Observe that for the values in question $\gamma=k\alpha,\ \delta=k\beta$,
we have $P = \alpha(1 + k x^2),\ Q = \beta(1 + k x^2)$; so that, writing
$P' = \alpha,\ Q' = \beta$, we have for $y$ the value
\[
\frac{1-y}{1+y} \;=\; \frac{(1-x)}{(1+x)}\,\frac{(P' - Q' x)^2}{(P' + Q' x)^2},
\]
which is the formula for a cubic transformation.

\medskip
\textbf{20.} It is important to notice that we cannot by writing
$\alpha=0$ or $\delta=0$ reduce the transformation to a quintic one; in
fact, the equation $k^2\alpha^2 = \Omega\delta^2$ shows that if either of
these equations is satisfied the other is also satisfied; and we have then
the foregoing case $\alpha=0,\ \delta=0$, giving not a quintic but a cubic
transformation.

And for the same reason we cannot by writing $\alpha=0,\ \beta=0,\ \gamma=0$
or $\beta=0,\ \gamma=0,\ \delta=0$ reduce the transformation to the order
$1$. There is thus no factor $\Omega - 1$.

\section*{[p.~125]}

\textbf{21.} As regards the non-existence of the factor $\Omega - 1$, I
further verify this by writing in the equations $\Omega = 1$; they thus
become
\[
k\alpha^2 = \delta^2,
\]
\[
k\,(2\alpha\gamma + 2\alpha\beta + \beta^2) = \gamma^2 + 2\gamma\delta + 2\beta\delta,
\]
\[
\gamma^2 + 2\beta\gamma + 2\alpha\delta + 2\beta\delta = k\,(2\alpha\gamma + \gamma^2 + 2\alpha\delta + \beta^2),
\]
\[
\delta^2 + 2\gamma\delta = k\,(\alpha^2 + 2\alpha\beta),
\]
which it is to be shown cannot be satisfied in general, but only for
certain values of $k$. Reducing the last equation, this is
$\gamma\delta = k\,\alpha\beta$, which, combined with the first, gives
$\alpha = \pm \sqrt{k}\,\alpha\gamma = \beta\delta$; and if for convenience
we assume $\alpha=1$, and write also $k = \theta^2$ (that is,
$k = \theta^2$), then the values of $\alpha,\beta,\gamma,\delta$ are
$\alpha=1,\ \beta = \gamma\theta^{-1},\ \gamma=\gamma,\ \delta=\theta$; which
values, substituted in the second and the third equations, give two
equations in $\gamma, \theta$; and from these, eliminating $\gamma$, we
obtain an equation for the determination of $\theta$, that is, of $k$. In
fact, the second equation gives
\[
\theta^2\,(2\gamma + 2\gamma\theta^{-1} + \gamma^2\theta^{-2})
\;=\; \gamma^2 + 2\gamma\theta + 2\gamma,
\]
or, dividing by $\gamma$ and reducing,
\[
\gamma\,(1 - \theta^2) = 2\theta^2(\theta^2 - 1)(\theta^2 - \theta + 1),
\]
that is,
\[
\gamma\,(1 + \theta^2) = -2\theta^2(\theta^2 - \theta + 1),
\]
or, as this may also be written,
\[
(\gamma + \theta^2)(1 + \theta^2) = -\theta^2(\theta - 1)^2,
\]
that is,
\[
\gamma + \theta^2 \;=\; \frac{-\theta^2(\theta - 1)^2}{\theta^2 + 1}.
\]

Moreover the third equation gives
\[
\gamma^2 + 2\gamma^2\theta^{-1} + 2\theta^2 + 2\gamma \;=\; \theta^2\,(2\gamma + 2\gamma^2\theta + 2\theta + \gamma^2\theta^{-1}),
\]
that is,
\[
\gamma^2\,(\theta^4 - 2\theta^3 + 2\theta - 1) - 2(\gamma + \theta^2)\,\theta^2(\theta^2 - 1) = 0;
\]
or dividing by $\theta^2 - 1$, it is
\[
\gamma^2\,(\theta - 1)^2 = 2\theta\,(\gamma + \theta^2);
\]
whence also
\[
\gamma \;=\; \frac{-2\theta}{(\theta - 1)^2}.
\]

Also
\[
4\theta^2(\theta^2 - \theta + 1)^2 = \gamma^2\,(\theta^2 + 1)^2,
\]
wherefore
\[
2(\theta^2 - \theta + 1)^2 = -\theta(\theta^2 + 1)
\quad\text{or}\quad 2(\theta^2 - \theta + 1)^2 + \theta(\theta^2 + 1) = 0,
\]
or
\[
\theta(\theta^2 + 1)^2 + 2(\theta^2 - \theta + 1)^2 = 0,
\]
that is,
\[
2\theta^4 - 3\theta^3 + 6\theta^2 - 2\theta + 2 = 0,
\]
or finally
\[
(2\theta^2 - \theta + 1)(\theta^2 - \theta + 2) = 0.
\]

\section*{[p.~126]}

We have thus $(2\theta + 1)^2 = \theta$, that is, $4\theta^2 + 3\theta + 1 = 0$
or $4k^2 + 3k + 1 = 0$, or else $(\theta^2 + 2)^2 = \theta$, that is,
$\theta^4 + 3\theta^2 + 4 = 0$ or $k^2 + 3k + 4 = 0$; viz.\ the equation in
$k$ is
\[
(4k^2 + 3k + 1)(k^2 + 3k + 4) = 0,
\]
these being in fact the values of $k$ given by the modular equation on
putting therein $\Omega = 1$.

The equation of the order $32$ thus contains the factor $\{(\Omega, 1)^4\}$
at least twice, and it does not contain either the factor $\Omega - 1$, or
the factor $\{(\Omega, 1)^6\}$ belonging to the quintic transformation; it
may be conjectured that the factor $\{(\Omega, 1)^4\}$ presents itself six
times, and that the form is
\[
\bigl\{(\Omega, 1)^4\bigr\}^6 \bigl\{(\Omega, 1)^8\bigr\} = 0;
\]
but I am not able to verify this, and I do not pursue the discussion
further.

\medskip
\textbf{22.} The foregoing considerations show the grounds of the
difficulty of the purely algebraical solution of the problem; the required
results, for instance the modular equation, are obtained not in the simple
form, but accompanied with special factors of high order. The
transcendental theory affords the means of obtaining the results in the
proper form without special factors; and I proceed to develop the theory,
reproducing the known results as to the modular and multiplier equations,
and extending it to the determination of the transformation-coefficients
$\alpha,\beta,\ldots$.

\section*{[p.~126] The Modular Equation. Art.\ Nos.\ 23 to 28.}

\textbf{23.} Writing, as usual, $q = e^{-\pi K'/K}$, we have $u$, a given
function of $q$, viz.
\[
u = \sqrt{2}\,q^{\frac{1}{8}}\,\frac{1+q\cdot 1+q^3\cdot 1+q^5\cdots}{1+q\cdot 1+q^2\cdot 1+q^3\cdots},
\]
\[
v = \sqrt{2}\,q^{\frac{1}{8}}\bigl(1 - q + 2q^2 - 3q^3 + 4q^4 - 6q^5 + 9q^6 - 14q^7 + \cdots\bigr)
\]
\[
= \sqrt{2}\,q^{\frac{1}{8}} f(q) \ \text{suppose,}
\]
and this being so, the several values of $v$ and of the other quantities in
question are all given in terms of $q$.

The case chiefly considered is that of $n$ an odd prime; and unless the
contrary is stated it is assumed that this is so. We have then $n+1$
transformations corresponding to the same number $n+1$ of values of $v$;
these may be distinguished as $v_0, v_1, v_2, \ldots, v_n$; viz.\ writing
$\alpha$ to denote an imaginary $n$th root of unity, we have
\[
v_0 = (-)^{\frac{n^2-1}{8}}\sqrt{2}\,q^{\frac{n}{8}}f(q^n),\quad
v_1 = \sqrt{2}\,(\alpha q^{\frac{1}{n}})^{\frac{1}{8}} f(\alpha q^{\frac{1}{n}}),
\]
\[
v_2 = \sqrt{2}\,(\alpha^2 q^{\frac{1}{n}})^{\frac{1}{8}} f(\alpha^2 q^{\frac{1}{n}}),\ \&c.,
\]
\[
v_n = \sqrt{2}\,q^{\frac{1}{8n}} f(q^{\frac{1}{n}}).
\]
(Observe $(-)^{\frac{n^2-1}{8}} = +1$ for $n = 8p\pm 1$, $= -1$ for $n = 8p\pm 3$.)

\section*{[p.~127]}

The occurrence of the fractional exponent $\tfrac{1}{8}$ is, as will appear,
a circumstance of great importance; and it will be convenient to introduce
the term ``octicity,'' viz.\ an expression of the form $q^{\frac{f}{8}}F(q)$
($f = 0$, or a positive integer not exceeding $7$, $F(q)$ a rational
function of $q$) may be said to be of the octicity $f$.

\medskip
\textbf{24.} The modular equation is of course
\[
(v - v_0)(v - v_1)\cdots(v - v_n) = 0;
\]
say this is
\[
v^{n+1} - A v^n + B v^{n-1} - \cdots = 0,
\]
so that $A = \Sigma v_0$, $B = \Sigma v_0 v_1$, \&c. In the development of these
expressions, the terms having a fractional exponent, with denominator $n$,
would disappear of themselves, as involving symmetrically the several
$n$th roots of unity; and each coefficient would be of the form
$q^{\frac{a}{8}} F(q)$, $F$ a rational and integral function of $q$. It is
moreover easy to see that, for the several coefficients $A, B, C, \ldots$,
$g$ will denote the positive residue (mod.\ 8) of $n, 2n, 3n,\ldots$
respectively.

Hence assuming, as the fact is, that these coefficients are severally
rational and integral functions of $q$, it follows that the form is
\[
a u^g + b u^{g+8} + c u^{g+16} + \cdots,
\]
$g$ having the foregoing values for the several coefficients respectively.
And it being known that the modular equation is as regards $u$ of the order
$=n+1$, there is a known limit to the number of terms in the several
coefficients respectively. We have thus for each coefficient an identity of
the form
\[
A = a u^g + b u^{g+8} + \cdots,
\]
$A$ and $u$ being each of them given in terms of $q$, the values of the
numerical coefficients $a, b, \ldots$ can be determined; and we thus arrive
at the modular equation.

\medskip
\textbf{25.} It is in effect in this manner that the modular equations are
calculated in Sohnke's Memoir. Various relations of symmetry in regard to
$(u, v)$ and other known properties of the modular equation are made use of
in order to reduce the number of the unknown coefficients to a minimum;
and (what in practice is obviously an important simplification) instead of
the coefficients $\Sigma v_0$, $\Sigma v_0 v_1$, \&c., it is the sums of
powers $\Sigma v_0,\ \Sigma v_0^2$, \&c., which are compared with their
expressions in terms of $u$, in order to the determination of the unknown
numerical coefficients $a, b, \ldots$. The process is a laborious one
(although less so than perhaps might beforehand have been imagined),
involving very high numbers; it requires the development up to high powers
of $q$, of the high powers of the before-mentioned function $f(q)$; and
Sohnke gives a valuable Table, which I reproduce here; adding to it the
three columns which relate to $\phi q$.

\clearpage
\section*{[p.~128] Sohnke's Table (with $\phi q$ columns added)}

\textit{[Table: Sohnke's Table of the developments in powers of $q$ of
$f^N(q)$ for $N = 3, 5, 7, 9, 11, 13$, and of $\phi q$ in three additional
columns. The table is here transcribed in landscape format; rows are
indexed by ind.\ of $q$ ($0, 1, 2, \ldots, 25$) and each column shows the
coefficient of the corresponding power of $q$ in the development of the
function indicated. Asterisked footnote: ``(Wrongly given by Sohnke as
$+34805616$).'']}

\begin{sidewaystable}
\renewcommand{\arraystretch}{0.9}
\centering
\footnotesize
\begin{tabular}{r|r|r|r|r|r|r|r|r|r|r|r|r}
\hline
ind.\ of $q$ & $\phi q$ & $\phi q^2$ & $\phi q^4$ & $f^3 q$ & $f^5 q$ & $f^7 q$ & $f^9 q$ & $f^{11} q$ & $f^{13} q$ & $f^{15} q$ & $f^{17} q$ & $f^{19} q$ \\
\hline
 0 & $+1$ & $+1$ & $+1$ & $+1$ & $+1$ & $+1$ & $+1$ & $+1$ & $+1$ & $+1$ & $+1$ & $+1$ \\
 1 & $+2$ & $0$ & $0$ & $-3$ & $-5$ & $-7$ & $-9$ & $-11$ & $-13$ & $-15$ & $-17$ & $-19$ \\
 2 & $+2$ & $+2$ & $0$ & $0$ & $+5$ & $+14$ & $+27$ & $+44$ & $+65$ & $+90$ & $+119$ & $+152$ \\
 3 & $0$ & $0$ & $0$ & $+5$ & $+5$ & $+7$ & $-9$ & $-55$ & $-130$ & $-238$ & $-380$ & $-552$ \\
 4 & $+2$ & $0$ & $+2$ & $0$ & $-15$ & $-49$ & $-90$ & $-110$ & $-65$ & $+90$ & $+391$ & $+874$ \\
 5 & $0$ & $0$ & $0$ & $-7$ & $+5$ & $+50$ & $+126$ & $+253$ & $+428$ & $+621$ & $+782$ & $+836$ \\
 6 & $+4$ & $+2$ & $0$ & $0$ & $+15$ & $+24$ & $-78$ & $-264$ & $-560$ & $-940$ & $-1356$ & $-1551$ \\
 7 & $0$ & $0$ & $0$ & $+3$ & $-30$ & $-141$ & $-322$ & $-440$ & $-260$ & $+450$ & $+1751$ & $+3704$ \\
 8 & $+2$ & $0$ & $0$ & $+15$ & $-10$ & $-49$ & $+108$ & $+803$ & $+2275$ & $+4860$ & $+8755$ & $+13889$ \\
 9 & $0$ & $0$ & $0$ & $-30$ & $+60$ & $+343$ & $+1056$ & $+2255$ & $+3978$ & $+5984$ & $+7565$ & $+8038$ \\
10 & $+4$ & $+2$ & $+2$ & $0$ & $+5$ & $+105$ & $+126$ & $-484$ & $-2860$ & $-7800$ & $-15691$ & $-26145$ \\
11 & $0$ & $0$ & $0$ & $+22$ & $-44$ & $-176$ & $-855$ & $-2783$ & $-6760$ & $-13325$ & $-22236$ & $-31958$ \\
12 & $+4$ & $0$ & $0$ & $-30$ & $-25$ & $+147$ & $+1437$ & $+5302$ & $+12870$ & $+24115$ & $+36490$ & $+44574$ \\
13 & $0$ & $0$ & $0$ & $-30$ & $+150$ & $+575$ & $+1485$ & $+2310$ & $+1820$ & $-2730$ & $-15470$ & $-42636$ \\
14 & $+4$ & $0$ & $0$ & $+10$ & $-25$ & $-630$ & $-2925$ & $-7700$ & $-15730$ & $-26280$ & $-36465$ & $-39738$ \\
15 & $0$ & $0$ & $0$ & $+90$ & $-126$ & $-651$ & $-1480$ & $+275$ & $+9100$ & $+30856$ & $+72475$ & $+139194$ \\
16 & $+4$ & $+2$ & $0$ & $-44$ & $+200$ & $+875$ & $+4995$ & $+14157$ & $+31525$ & $+59565$ & $+95485$ & $+126465$ \\
17 & $0$ & $0$ & $0$ & $-77$ & $+275$ & $+1729$ & $+6390$ & $+18887$ & $+45045$ & $+91200$ & $+158015$ & $+233358$ \\
18 & $+4$ & $0$ & $+2$ & $+0$ & $-225$ & $-1407$ & $-7479$ & $-26015$ & $-67925$ & $-149600$ & $-282115$ & $-451250$ \\
19 & $0$ & $0$ & $0$ & $+132$ & $-220$ & $-1995$ & $-9504$ & $-30855$ & $-77220$ & $-159510$ & $-280440$ & $-422256$ \\
20 & $+6$ & $0$ & $0$ & $-110$ & $+250$ & $+2128$ & $+10934$ & $+39226$ & $+109110$ & $+253935$ & $+512435$ & $+918680$ \\
21 & $0$ & $0$ & $0$ & $-115$ & $+675$ & $+2870$ & $+10560$ & $+27357$ & $+50500$ & $+59450$ & $+30450$ & $-67450$ \\
22 & $+4$ & $+2$ & $0$ & $-160$ & $+0$ & $+1638$ & $+10395$ & $+39446$ & $+114550$ & $+277200$ & $+583240$ & $+1097290$ \\
23 & $0$ & $0$ & $0$ & $+198$ & $-825$ & $-3990$ & $-16236$ & $-58300$ & $-176580$ & $-449820$ & $-996300$ & $-1982880$ \\
24 & $+4$ & $0$ & $0$ & $+165$ & $-825$ & $-4500$ & $-18810$ & $-69685$ & $-216905$ & $-589580$ & $-1387845$ & $-2901620$ \\
25 & $0$ & $0$ & $0$ & $-280$ & $+475$ & $+5390$ & $+25025$ & $+88825$ & $+255300$ & $+656535$ & $+1494125$ & $+3060180$ \\
\hline
\end{tabular}

\vspace{1ex}
{\footnotesize $\asterism$ Wrongly given by Sohnke as $+34805616$.}
\end{sidewaystable}

\clearpage
\section*{[p.~130] Series of modular equations (after Sohnke).}

\textbf{26.} I give from Sohnke the series of modular equations, adding
those for the composite cases $n=9$ and $n=15$, as to which see the
remarks which follow the Table.

\medskip
\textit{[Table: modular equation for $n=3$. The grid has rows
$u^4, u^3, u^2, u, 1$ and columns $v^4, v^3, v^2, v, 1$; the non-zero
entries are]}
\[
\begin{aligned}
&u^4: \quad -1 \ \text{at column }1, \\
&u^3: \quad +2 \ \text{at column }v^2, \\
&u^2: \quad 0, \\
&u^1: \quad -2 \ \text{at column }v^2, \\
&1: \quad +1 \ \text{at column }v^4.
\end{aligned}
\]
The column-sum reads
\[
1 \cdot v^4 + 2\cdot v^2 \cdot u^3 + 0 - 2\cdot v^2 u - 1\cdot u^4 = (v+1)^3 (v-1).
\]

\medskip
\textit{[Table: modular equation for $n=5$. Rows $u^6, u^5, u^4, u^3, u^2, u, 1$,
columns $v^6, v^5, v^4, v^3, v^2, v, 1$; the non-zero entries are]}
\[
\begin{aligned}
&u^6: \quad -1 \ \text{at column }1, \\
&u^5: \quad +4 \ \text{at column }v^2, \\
&u^4: \quad -5 \ \text{at column }v^4, \\
&u^3: \quad +5 \ \text{at column }v^3, \\
&u^2: \quad -4 \ \text{at column }v^5, \\
&u: \ \ \quad +1 \ \text{at column }v^6.
\end{aligned}
\]
The full equation reads
\[
1 + 4 + 5 + 0 - 5 - 4 - 1 = (v+1)^5 (v-1).
\]

\medskip
\textit{[Table: modular equation for $n=7$. Rows $u^8,\ldots,u,1$, columns
$v^8,\ldots,v,1$; the non-zero entries form a skew-diagonal pattern with
binomial-like values:]}
\[
\begin{aligned}
&u^8: \quad +1 \ \text{at column }1, \\
&u^7: \quad -8 \ \text{at column }v^2, \\
&u^6: \quad +28 \ \text{at column }v^4, \\
&u^5: \quad -56 \ \text{at column }v^5, \\
&u^4: \quad +70 \ \text{at column }v^4, \\
&u^3: \quad -56 \ \text{at column }v^5, \\
&u^2: \quad +28 \ \text{at column }v^6, \\
&u: \ \ \quad -8 \ \text{at column }v^7, \\
&1: \ \ \quad +1 \ \text{at column }v^8.
\end{aligned}
\]
The complete column-sum reads
\[
+1 - 8 + 28 - 56 + 70 - 56 + 28 - 8 + 1 = (v-1)^8.
\]

\section*{[p.~131]}

\textit{[Table: modular equation for $n=9$. Rows $u^{12}, u^{11}, \ldots, u, 1$
and columns $v^{12}, v^{11}, \ldots, v, 1$. The diagonal pattern of
non-zero entries is]}
\[
\begin{aligned}
&u^{12}: \ +1 \ \text{at col }1, \\
&u^{11}: \ -8 \ \text{at col }v^3, \\
&u^{10}: \ +16 \ \text{at col }v^6, \\
&u^9: \ -16 \ \text{at col }v^8,\quad +10 \ \text{at col }v^4, \\
&u^8: \ +8 \ \text{at col }v^7,\quad -16 \ \text{at col }v^5, \\
&u^7: \ -16 \ \text{at col }v^6,\quad +15 \ \text{at col }v^4, \\
&u^6: \ +48 \ \text{at col }v^5,\quad -84 \ \text{at col }v^3, \\
&u^5: \ -84 \ \text{at col }v^4,\quad +48 \ \text{at col }v^6, \\
&u^4: \ +48 \ \text{at col }v^3,\quad +15 \ \text{at col }v^5, \\
&u^3: \ +15 \ \text{at col }v^4,\quad -24 \ \text{at col }v^5, \\
&u^2: \ -24 \ \text{at col }v^5,\quad +10 \ \text{at col }v^3,\quad -16 \ \text{at col }v^4, \\
&u: \ \ +10 \ \text{at col }v^4,\quad +16 \ \text{at col }v^2,\quad -16 \ \text{at col }v^3, \\
&1: \ \ +8 \ \text{at col }v^3,\quad -16 \ \text{at col }v^4,\quad +1 \ \text{at col }v.
\end{aligned}
\]
The full column-sum, when read across $v^{12}, v^{11}, \ldots, v, 1$, is
\[
1 - 8 + 26 - 40 + 15 + 48 - 84 + 48 + 15 - 40 + 26 - 8 + 1
\;=\; (v - 1)^{10}(v + 1)^2.
\]

\medskip
\textit{[Table: modular equation for $n=11$. Rows $u^{12},u^{11},\ldots,u,1$
and columns $v^{12},v^{11},\ldots,v,1$. Non-zero entries:]}
\[
\begin{aligned}
&u^{12}: \ -1 \ \text{at col }1, \\
&u^{11}: \ +32 \ \text{at col }v^2, \\
&u^{10}: \ -44 \ \text{at col }v^4, \\
&u^9: \ +88 \ \text{at col }v^3,\quad +22 \ \text{at col }v^7, \\
&u^8: \ -165 \ \text{at col }v^5, \\
&u^7: \ +132 \ \text{at col }v^4, \\
&u^6: \ +44 \ \text{at col }v^5,\quad -44 \ \text{at col }v^8, \\
&u^5: \ -132 \ \text{at col }v^6, \\
&u^4: \ +165 \ \text{at col }v^7, \\
&u^3: \ -22 \ \text{at col }v^8,\quad -88 \ \text{at col }v^9, \\
&u^2: \ +44 \ \text{at col }v^9, \\
&u: \ \ +22 \ \text{at col }v^{10},\quad -32 \ \text{at col }v^{11}, \\
&1: \ \ +1 \ \text{at col }v^{12}.
\end{aligned}
\]
The full row of column-sums reads
\[
1 + 10 + 44 + 110 + 165 + 132 + 0 - 132 - 165 - 110 - 44 - 10 - 1 \;=\; (v+1)^{11}(v - 1).
\]

\section*{[p.~132]}

\textit{[Table: modular equation for $n=13$. Rows $u^{14},u^{13},\ldots,u,1$
and columns $v^{14},v^{13},\ldots,v,1$. Non-zero entries are arranged in
the skew-diagonal pattern characteristic of $n \equiv 5\pmod{8}$:]}
\[
\begin{aligned}
&u^{14}: \ -1 \ \text{at col }1, \\
&u^{13}: \ +64 \ \text{at col }v^2,\quad -52 \ \text{at col }v^8, \\
&u^{12}: \ -65 \ \text{at col }v^4, \\
&u^{11}: \ +208 \ \text{at col }v^6, \\
&u^{10}: \ +208 \ \text{at col }v^5,\quad -429 \ \text{at col }v^7, \\
&u^9: \ +520 \ \text{at col }v^6,\quad +52 \ \text{at col }v^9, \\
&u^8: \ -429 \ \text{at col }v^7, \\
&u^7: \ -208 \ \text{at col }v^7, \\
&u^6: \ +429 \ \text{at col }v^8, \\
&u^5: \ -52 \ \text{at col }v^{11},\quad -520 \ \text{at col }v^9, \\
&u^4: \ +429 \ \text{at col }v^{10}, \\
&u^3: \ -208 \ \text{at col }v^{11}, \\
&u^2: \ +65 \ \text{at col }v^{12}, \\
&u: \ \ +52 \ \text{at col }v^{13},\quad -64 \ \text{at col }v^{13}, \\
&1: \ \ +1 \ \text{at col }v^{14}.
\end{aligned}
\]
The complete column-sum reads
\[
1 + 12 + 65 + 208 + 429 + 572 + 429 + 0 - 429 - 572 - 429 - 208 - 65 - 12 - 1 \;=\; (v+1)^{13}(v - 1).
\]

\clearpage
\section*{[p.~133] Modular equation for $n=17$.}

\textit{[Table: a large rotated table (sidewaystable in the original) giving
the modular equation for $n=17$. Rows are powers $u^{18}, u^{17}, \ldots, u, 1$
and columns are powers $v^{18}, v^{17}, \ldots, v, 1$. The entries follow
the skew-symmetric pattern of $n \equiv 1 \pmod{8}$ (perfect symmetry, no
sign reversal); the non-zero entries appearing in the printed table, read
along the principal diagonals, are coefficients such as $+1, -8, +28, +66,
+220, \ldots$ on the dexter (descending) diagonal and their reflections on
the sinister (rising) diagonal. The column-sum reduces to
$(v - 1)^{18}$.]}

\clearpage
\section*{[p.~134] Modular equation for $n=19$.}

\textit{[Table: rotated landscape table for the modular equation of
$n=19$. Rows are $u^{20}, u^{19}, \ldots, u, 1$ and columns are
$v^{20}, v^{19}, \ldots, v, 1$. Owing to the case $n \equiv 3 \pmod{8}$, the
table is not perfectly symmetric but has alternating sign reversals on the
sinister diagonal. The non-zero diagonal coefficients begin
$-1,\ +114,\ -114,\ 0,\ -2584,\ -6859,\ -2584,\ldots$ as referred to in
Art.\ 26. The full column-sum reduces to $(v + 1)^{19}(v - 1)$.]}

\clearpage
\section*{[p.~135] Modular equation for $n=21$ and remarks.}

\textit{[Table: rotated landscape table for the modular equation of
$n=21$. Rows are $u^{22}, u^{21}, \ldots, u, 1$ and columns are
$v^{22}, v^{21}, \ldots, v, 1$. The pattern is symmetric with sign
reversals on alternate lines parallel to the sinister diagonal. The
column-sum reduces to $(v + 1)^{21}(v - 1)$.]}

\section*{[p.~136] Remarks on the Tables.}

Various remarks arise on the Tables. Attending first to the cases $n$ a
prime number; the only terms of the order $n+1$ in $v$ or $u$ are
$v^{n+1} \mp u^{n+1}$, viz.\ $n = 3$ or $5$ (mod.\ 8) the sign is $-$, but
$n = 1$ or $7$ (mod.\ 8) the sign is $+$. And there is in every case a
pair of terms $v^n u^a$ and $v^a u^n$, having coefficients equal in absolute
magnitude, but of opposite signs, or of the same sign, in the two cases
respectively.

Each Table is symmetrical in regard to its two diagonals respectively, so
that every non-diagonal coefficient occurs (with or without reversal of
sign) $4$ times; viz.\ in the case $n=1$ or $7$ (mod.\ 8) this is a
perfect symmetry, without reversal of sign; but in the case $n=3$ or $5$
(mod.\ 8) it is, as regards the lines parallel to either diagonal, and in
regard to the other diagonal, alternately a perfect symmetry without
reversal of sign and a skew symmetry with reversal. Thus in the case
$n = 19$, the lines parallel to the dexter diagonal are
$-1$ (symmetrical), $+114, -114$ (skew), $0, -2584, -6859, -2584,$ \ldots
(symmetrical), and so on. The same relation of symmetry is seen in the
composite cases $n=9$ and $n=15$, both belonging to $n=1$ or $7$, mod.\ 8.

If, as before, $n$ is prime, then putting in the modular equation $u=1$,
the equation in the case $n=1$ or $7$ (mod.\ 8) becomes $(v - 1)^{n+1} = 0$,
but in the case $n=3$ or $5$ (mod.\ 8) it becomes $(v + 1)^n (v - 1) = 0$.

\medskip
\textbf{27.} In the case $n$ a composite number not containing any square
factor, then dividing $n$ in every possible way into two factors $n = a b$
(including the divisions $n\cdot 1$ and $1\cdot n$), and denoting by
$\beta$ an imaginary $b$th root of unity, a value of $v$ is
\[
\pm \sqrt{2}\,(\beta q^{\frac{a}{b}})^{\!\frac{1}{8}}\, f\!\left(\beta q^{\frac{a}{b}}\right);
\]
so that the whole number of roots (or order of the modular equation) is
$=\nu$, if $\nu$ be the sum of the divisors of $n$. Thus $n=15$, the values
are
\[
\sqrt{2}\,q^{\frac{15}{8}}f(q^{15}),\quad
-\sqrt{2}\,q^{\frac{5}{8}}f(q^{5}),\quad
-\sqrt{2}\,q^{\frac{3}{8}}f(q^3),\quad
\sqrt{2}\,q^{\frac{1}{8}}f(q),
\]
\[
1\ ,\quad 3\ ,\quad 5\ ,\quad 15\ \text{roots};
\]
and the order of the modular equation is $=24$. The modular equation
might thus be obtained as for a prime number; but it is easier to
decompose $n$ into its prime factors, and consider the transformation as
compounded of transformations of these prime orders. Thus $n = 15$, the
transformation is compounded of a cubic and a quintic one. If the $v$ of
the cubic transformation be denoted by $\theta$, then we have
\[
\theta^4 + 2\theta^3 u^3 - 2\theta u - u^4 = 0;
\]
and to each of the four values of $\theta$ corresponds the six values of $v$
belonging to the quintic transformation given by
\[
\ldots
\]

The equation in $v$ is thus
\[
\ldots
\]

\section*{[p.~137]}

where $\theta_1, \theta_2, \theta_3, \theta_4$ are the roots of the equation in $\theta$, viz.\ we have
\[
\theta^4 + 2\theta^3 u^3 - 2\theta u - u^4 \;=\; (\theta - \theta_1)(\theta - \theta_2)(\theta - \theta_3)(\theta - \theta_4);
\]
and it was in this way that the equation for the case $n=15$ was calculated.
Observe that writing $u=1$, we have $(\theta + 1)^3 (\theta - 1) = 0$, or say
$\theta_1 = \theta_2 = \theta_3 = -1,\ \theta_4 = +1$. The equation in $v$
thus becomes $\{(v - 1)^3(v + 1)\}^3 (v - 1)^3(v + 1) = 0$, that is,
$(v - 1)^{12}(v + 1)^{12} = 0$.

\medskip
\textbf{28.} The case where $n$ has a square factor is a little different;
thus $n=9$, the values are
\[
1\ ,\quad 3\ ,\quad 9\ \ldots\ \text{roots},
\]
but here $\omega$ being an imaginary cube root of unity, the second term
denotes the three values,
\[
\sqrt{2}\,q^{\frac{1}{8}}f(q),\quad
\sqrt{2}\,(\omega q)^{\frac{1}{8}}f(\omega q),\quad
\sqrt{2}\,(\omega^2 q)^{\frac{1}{8}}f(\omega^2 q),
\]
the first of which is $=u$, and is to be rejected; there remain $1+2+9$,
$=12$ roots, or the equation is of the order $12$.

Considering the equation as compounded of two cubic transformations, if the
value of $v$ for the first of these be $\theta$, then we have
\[
\theta^4 + 2\theta^3 u^3 - 2\theta u - u^4 = 0;
\]
and to the four values of $\theta$ correspond severally the four values of
$v$ given by the equation
\[
v^4 + 2 v^3 \theta^3 - 2 v \theta - \theta^4 = 0.
\]

One of these values is however $v = -u$, since the $u^4$-equation is
satisfied on writing therein $v = -u$; hence, writing
\[
\theta^4 + 2\theta^3 u^3 - 2\theta u - u^4 \;=\; (\theta - \theta_1)(\theta - \theta_2)(\theta - \theta_3)(\theta - \theta_4),
\]
we have an equation
\[
(v^4 + 2 v^3 \theta_1^3 - 2 v \theta_1 - \theta_1^4)(v^4 + \cdots\theta_2^3\cdots\theta_2^4)\cdots
(v^4 + \cdots\theta_4^3\cdots\theta_4^4) = 0,
\]
which contains the factor $(v + u)^4$ and, divested hereof, gives the
required modular equation of the order $12$; it was in fact obtained in
this manner.

Observe that writing $u=1$, we have $(\theta + 1)^3(\theta - 1) = 0$, or say
$\theta_1 = \theta_2 = \theta_3 = -1,\ \theta_4 = 1$;

the modular equation then becomes
\[
\{(v - 1)^3(v + 1)\}^3\,(v + 1)^3\,(v - 1)^3(v + 1)^3 = 0,
\]
that is,
\[
(v - 1)^9(v + 1)^9 = 0.
\]

\clearpage
\section*{[p.~138] The Multiplier Equation. Art.\ No.\ 29.}

\textbf{29.} The theory is in many respects analogous to that of the
modular equation. To each value of $v$ there corresponds a single value of
$M$; hence $M$, or what is the same thing $\dfrac{1}{M}$, is determined by
an equation of the same order as $v$, viz.\ $n$ being prime, the order is
$=n+1$. The last term of the equation is constant, and the other
coefficients are rational and integral functions of $u^4$, of a degree not
exceeding $\tfrac{1}{2}(n-1)$; and not only so, but they are,
$n \equiv 1$ (mod.\ 4), rational and integral functions of $u^4(1 - u^4)$,
and $n \equiv 3$ (mod.\ 4), alternately of this form and of the same form
multiplied by the factor $(1 - 2u^4)$.

The values are in fact given as transcendental functions of $q$; viz.\
denoting by $M_0, M_1, M_2, \ldots, M_n$ the values corresponding to
$v_0, v_1, v_2, \ldots, v_n$ respectively, and writing
\[
\phi(q) = \frac{(1+q)(1+q^2)(1+q^3)\cdots\,(1 - q^2)(1-q^4)\cdots}
              {(1-q)(1-q^3)(1-q^5)\cdots\,(1+q^2)(1+q^4)\cdots},
\]
\[
\quad\quad\quad = 1 + 2q + 2q^4 + 2q^9 + 2q^{16} + \cdots,
\]
then we have
\[
(-)^{\frac{n-1}{2}} M_0^{\frac{1}{2}} \;=\; \frac{\phi(q)}{\phi(q^n)},
\]
\[
M_1^{\frac{1}{2}} \;=\; \frac{\phi(q)}{\phi(\alpha q^{\frac{1}{n}})},
\quad\text{($\alpha$ an imaginary $n$th root of unity)},
\]
\[
M_n^{\frac{1}{2}} \;=\; \frac{\phi(q)}{\phi(q^{\frac{1}{n}})}.
\]
Hence, the form of the equation being known, the values of the numerical
coefficients may be calculated; and it was in this way that Joubert
obtained the following results. I have in some cases changed the sign of
Joubert's multiplier, so that in every case the value corresponding to
$u=0$ shall be $M=1$.

The equations are:
\[
n = 3,\quad
\frac{1}{M^4} - 3 \;=\; 0; \qquad u = 0,\ \text{this is}\quad \left(\frac{1}{M} - 1\right)^{\!3}\!\left(\frac{1}{M} + 3\right) = 0;
\]
\[
\qquad\qquad u = 1,\ \text{it is}\quad \left(\frac{1}{M} + 1\right)^{\!3}\!\left(\frac{1}{M} - 3\right) = 0.
\]

\section*{[p.~139]}

\[
n = 5,\quad \frac{1}{M^6}\,,
\]
\[
+ \frac{1}{M^4}\cdot(-10),
\]
\[
+ \frac{1}{M^3}\cdot(+35),
\]
\[
+ \frac{1}{M^2}\cdot(-60),
\]
\[
+ \frac{1}{M}\cdot(+55),
\]
\[
\frac{1}{M}\cdot\bigl[-26 + 256\,u^4(1 - u^4)\bigr] + 5 = 0;
\]
\[
u = 0\ \text{or}\ 1,\ \text{this is}\quad \left(\frac{1}{M} - 1\right)^{\!5}\!\left(\frac{1}{M} - 5\right) = 0.
\]

\[
n = 7,\quad \frac{1}{M^8},
\]
\[
+ \frac{1}{M^7}\cdot 0,
\]
\[
+ \frac{1}{M^6}\cdot(-28),
\]
\[
+ \frac{1}{M^5}\cdot\bigl[+112(1 - 2u^4)\bigr],
\]
\[
+ \frac{1}{M^4}\cdot(-210),
\]
\[
+ \frac{1}{M^3}\cdot\bigl[+224(1 - 2u^4)\bigr],
\]
\[
+ \frac{1}{M^2}\cdot\bigl[-140 - 21\cdot 25\,6\,u^4(1 - u^4)\bigr],
\]
\[
+ \frac{1}{M}\cdot\bigl[48 + 2048\,u^4(1 - u^4)\bigr](1 - 2u^4) + 7 = 0;
\]
\[
u = 0,\ \text{this is}\quad \left(\frac{1}{M} - 1\right)^{\!7}\!\left(\frac{1}{M} - 7\right) = 0;
\]
\[
u = 1,\ \text{it is}\quad \left(\frac{1}{M} + 1\right)^{\!7}\!\left(\frac{1}{M} + 7\right) = 0.
\]

\[
n = 11,\quad \frac{1}{M^{12}},
\]
\[
+ \frac{1}{M^{11}}\cdot 0,
\]
\[
+ \frac{1}{M^{10}}\cdot(-66),
\]
\[
+ \frac{1}{M^9}\cdot\bigl[+440(1 - 2u^4)\bigr],
\]

\section*{[p.~140]}

\[
+ \frac{1}{M^8}\cdot(-1485),
\]
\[
+ \frac{1}{M^7}\cdot\bigl[+3168(1 - 2u^4)\bigr],
\]
\[
+ \frac{1}{M^6}\cdot\bigl[-4620 - 3\cdot 11^2\cdot 256\,u^4(1 - u^4)\bigr],
\]
\[
+ \frac{1}{M^5}\cdot\bigl[+4752 + 11\cdot 4096\,u^4(1 - u^4)\bigr]\,(1 - 2u^4),
\]
\[
+ \frac{1}{M^4}\cdot\bigl[-3465 - 3\cdot 7\cdot 11\cdot 512\,u^4(1 - u^4)\bigr],
\]
\[
+ \frac{1}{M^3}\cdot\bigl[+1760 - 11\cdot 83\cdot 2048\,u^4(1 - u^4)\bigr]\,(1 - 2u^4),
\]
\[
+ \frac{1}{M^2}\cdot\bigl[-594 - 9\cdot 11\cdot 37\cdot 256\,u^4(1 - u^4) - 3\cdot 11\cdot 131072\,[u^4(1 - u^4)]^2\bigr],
\]
\[
+ \frac{1}{M}\cdot\bigl[120 + 15\cdot 4096\,u^4(1 - u^4) - 524288\,(u^4(1 - u^4))^2\bigr]\,(1 - 2u^4)
\]
\[
\qquad - 11 = 0.
\]
\[
u = 1,\ \text{it is}\quad \left(\frac{1}{M} + 1\right)^{\!11}\!\left(\frac{1}{M} - 11\right) = 0.
\]

\section*{[p.~140] The Multiplier as a rational function of $u, v$. Art.\ Nos.\ 30 to 36.}

\textbf{30.} The multiplier $M$, as having a single value corresponding to
each value of $v$, is necessarily a rational function of $u, v$; and such
an expression of $M$ can, as remarked by K\"onigsberger, be deduced from
the multiplier equation by means of Jacobi's theorem,
\[
M^n \;=\; \frac{1}{n}\,\frac{\lambda(1 - \lambda^2)}{k\,(1 - k^2)}\,\frac{dk}{d\lambda},
\]
viz.\ substituting for $k, \lambda$ their values $u^4, v^4$, and observing
that if the modular equation be $F(u, v) = 0$ so that the value of
$\dfrac{dv}{du}$ is $-F'(v) \div F'(u)$, this is
\[
M^n \;=\; \frac{1}{n}\,\frac{(1 - v^4)\,v\, F'v}{(1 - u^4)\,u\, F'u};
\]
and then in the multiplier equation separating the terms which contain the
odd and even powers, and writing it in the form
$\Phi(M^2) + M\Psi(M^2) = 0$, this equation, substituting therein for $M^2$
its value, gives the value of $M$ rationally.

The rational expression of $M$ in terms of $u, v$ is of course
indeterminate, since its form may be modified in any manner by means of the
equation $F(u, v) = 0$; and in the expression obtained as above, the orders
of the numerator and the denominator are far too high. A different form
may be obtained as follows: for greater convenience I seek for the value
not of $M$ but of $\dfrac{1}{M}$.

\section*{[p.~141]}

\textbf{31.} Denoting, as above, by $M_0, M_1, \ldots, M_n$ the values
which correspond to $v_0, v_1, \ldots, v_n$ respectively, and writing
$S_{\frac{1}{M}} = \dfrac{1}{M_0} + \dfrac{1}{M_1} + \cdots + \dfrac{1}{M_n}$,
\&c., we have $S_{\frac{1}{M}}, S_{\frac{v}{M}}$, \&c., all of them
expressible as determinate functions of $u$; and we have moreover the
theorem that each of these is a rational and integral function of $u$: we
have thus the series of equations
\[
S_{\frac{1}{M}} = A,\quad S_{\frac{v}{M}} = B,\ \ldots,\ S_{\frac{v^n}{M}} = H,
\]
where $A, B, \ldots, H$ are rational and integral functions of $u$. These
give linearly the different values of $\dfrac{1}{M}$; in fact, we have
\[
(v_0 - v_1)\cdots(v_0 - v_n)\,\frac{1}{M_0}
\;=\; H - G\,\Sigma v_1 + F\,\Sigma v_1 v_2 - \cdots \pm A\,v_1 v_2 \cdots v_n,
\]
where $\Sigma v_1, \Sigma v_1 v_2$, \&c.\ denote the combinations formed
with the roots $v_1, v_2, \ldots, v_n$ (these can be expressed in terms of
the single root $v_0$); and we have also
$(v_0 - v_1)\cdots(v_0 - v_n) = F'(v_0)$ the resulting equation is
consequently
\[
F' v_0\,\frac{1}{M_0} \;=\; R(u, v_0),
\]
a determinate rational and integral function of $(u, v_0)$; but as the
same formula exists for each root of the modular equation, we may herein
write $M, v$ in place of $M_0, v_0$; and the formula thus is
\[
F'v\cdot\frac{1}{M} \;=\; R(u, v),
\]
viz.\ we thus obtain the required value of $\dfrac{1}{M}$ as a rational
function, the denominator being the determinate function $F'v$, and the
numerator being, as is easy to see, a determinate function of the order
$n$ as regards $v$.

\medskip
\textbf{32.} The method is applicable when $M$ is only known by its
expression in terms of $q$; but if we know for $M$ an expression in terms
of $v, u$, then the method transforms this into a standard form as above.
By way of illustration I will consider the case $n = 3$, where the modular
equation is
\[
v^4 + 2 v^3 u^3 - 2 v u - u^4 = 0,
\]
\[
v = u + \frac{1 - 2u^4}{\ldots},
\]
and where a known expression of $M$ is $\dfrac{1}{\ldots}$. Here writing
$S_{-1}, S_0\ (=4), S_1$, \&c.\ to denote the sum of the powers $-1, 0, 1$,
\&c.\ of the roots of the equation, we have
\[
S_{\frac{1}{M}} \;=\; 0 \;=\; S_3 + 2u^3 S_{-1},\quad\text{as appears from the values presently given,}
\]
\[
S_{\frac{v}{M}} \;=\; 0 \;=\; S_4 + 2u^3 S_0,
\]
\[
S_{\frac{v^2}{M}} \;=\; S_5 + 2u^3 S_1 \;=\; 6u;
\]

\section*{[p.~142]}

and observing that $v_0$, being ultimately replaced by $v$, we have
\[
\Sigma v_1 = \Sigma v_0 - v,\quad
\Sigma v_1 v_2 = \Sigma v_0 v_1 - v\Sigma v_0 + v^2,\quad
v_1 v_2 v_3 = \Sigma v_0 v_1 v_2 - v\Sigma v_0 v_1 + v^2\Sigma v_0 - v^3,
\]
that is,
\[
\Sigma v_1 = -2u^3 - v,\quad
\Sigma v_1 v_2 = -2u^3 v - v^2,\quad
v_1 v_2 v_3 = 2u - 2u^3 v^2 - v^3,
\]
we have
\[
F'v\cdot\frac{1}{M} \;=\; (S_3 + 2u^3 S_0)\ \ldots\ + (2u^3 + v)(S_2 + 2u^3 S_1)
\]
\[
\qquad + (2u^3 v + v^2)(S_1 + 2u^3 S_0)
\]
\[
\qquad + (-2u + 2u^3 v^2 + v^3)(S_0 + 2u^3 S_{-1}),
\]
viz.\ this is
\[
\ldots
\]

But we have
\[
S_{-1} = -3,\quad S_0 = 4,\quad S_1 = -2u^3,\quad S_2 = u^2,\quad S_3 = 6u - 8u^3,\ldots
\]
and the equation thus is
\[
(2v^3 + 3v^2 u^3 - u)\cdot\frac{1}{M} \;=\; 3\,(v^2 u^2 + 2u^3 v + 1)\,;
\]
to verify which observe that, substituting herein for $\dfrac{1}{M}$ its
value $1 + \dfrac{2u^3}{v}$, the equation becomes
\[
S v^3 u^3 + (2v^2 + 0)(v^2 u^3) S v u\cdot (v^2 u^2 + 2u^3 v + 1) \,;
\]
that is,
\[
2v^4 + 4u^3 v^3 - 4 v u - 2u^4 = 0,
\]
as it should do.

\medskip
\textbf{33.} Any expression whatever of $M$ in terms of $u, v$ is in fact
one of a system of four expressions; viz.\ we may simultaneously change
\[
\text{that is, signs are}
\]
\begin{center}
\begin{tabular}{c|c|cccc}
$u$ & $v$ & $\dfrac{1}{M}$ & $n=1$ & $n=3$ & $n=5$,\ $n=7\pmod 8$ \\[2pt]
\hline
$u$ & $v$ & $\dfrac{1}{M}$ & $-+++$ & $+-+-$ & $+--+$ \\[2pt]
into $v$ & $(-)^{\frac{n^2-1}{8}}u$ & $(-)^{\frac{n-1}{2}}\dfrac{n^2}{M^n}$ & $++++$ & $+++-$ & $+++-$ \\[2pt]
or $\dfrac{1}{u}$ & $\dfrac{1}{v}$ & $\dfrac{v^n}{u^n M}$ & $++++$ & $+++-$ & $+++-$ \\[2pt]
or $\dfrac{1}{v}$ & $(-)^{\frac{n^2-1}{8}}\dfrac{1}{u}$ & $(-)^{\frac{n-1}{2}}\dfrac{u^n}{v^n}\dfrac{n^2}{M^n}$ & $+++$ & $+-+$ & $+-+$ \\
\end{tabular}
\end{center}

\end{document}
